\section{Experimental setup}\label{sec:setup}

\subsection{Analysis plans}\label{sec:plans}

The design, estimands and tests were specified in advance in a version-controlled
repository, and Table~\ref{tab:chronology} in \ref{app:deviations} gives the chronology.
For the main study, the initial plan and its first amendment preceded the generation of
the evaluation candidates. The final amendment came after we had inspected
generator-side statistics (accuracy, share of extractable answers, truncation) and
before any verifier scored an evaluation problem. The plans of the two further studies, H6 and H7, with their problem manifests and
analysis code, were fixed before their
candidates were generated. Because we did not use an external registry, the commit
history is the timestamped record of each plan. The hypotheses are labelled c1--c5 in
the main study, H6a/b in H6 and H7a/b in H7. Every result outside these
tests is exploratory and labelled as such.

\subsection{Problems and candidate pools}\label{sec:pools}

Table~\ref{tab:populations} lists the five evaluation populations, and each row is
used as follows. \emph{Main, full} is the main evaluation set of 2,000 problems: 800 from
the GSM8K test split \cite{cobbe2021gsm8k}, 400 sampled from MATH500
\cite{lightman2023verify}, and 800 further MATH test problems \cite{hendrycks2021math}
outside MATH500, sampled with a fixed seed. It is used for the memory comparison c2 and
for the exploratory grid. \emph{Main, original manifest} is its first 1,200 problems
(GSM8K and MATH500). This subset was frozen before the 800 extra MATH problems were added,
and it is the prespecified population of c1. \emph{Main, MATH} contains the 1,200 MATH
problems of the main set (MATH500 and the extension) and is used for c4 and c5.
\emph{H6} (800 MATH Level-5 test problems) and \emph{H7} (1,060 MATH Level-4 and
Level-5 test problems) are the populations of the two further studies; none of their
problems appears anywhere else in this work. A separate development set of 400 problems
(200 GSM8K train, 200 MATH train) was used only for the go/no-go gate and for the choices
assigned to development data.

\begin{table}[!htb]
\centering
\small
\caption{Principal evaluation populations. ``New'' problems appear nowhere else in this
work (text-level exclusion against all other items and against all of
MATH500).}\label{tab:populations}
\resizebox{\linewidth}{!}{%
\begin{tabular}{lrll}
\toprule
population & $n$ & composition & used for \\
\midrule
Main, full & 2,000 & 800 GSM8K + 400 MATH500 + 800 MATH test & c2, Table~\ref{tab:main} \\
Main, original manifest & 1,200 & 800 GSM8K + 400 MATH500 & c1 \\
Main, MATH & 1,200 & 400 MATH500 + 800 MATH test & c4, c5 \\
H6 & 800 & MATH test Level 5 (new) & H6a/b \\
H7 & 1,060 & MATH test Level 4 (905) + Level 5 (155), new & H7a/b \\
\bottomrule
\end{tabular}}
\end{table}

The main generator is Qwen2.5-1.5B-Instruct \cite{qwen2.5}. For every problem it
produces a pool of 32 candidates (4 seeds $\times$ 8 samples, temperature 0.8,
top-$p$ 1.0, at most 1,024 new tokens) with the system prompt ``Please reason step by
step, and put your final answer within \texttt{\textbackslash boxed\{\}}.'' The pools
are generated once in BF16 with vLLM 0.30 \cite{kwon2023vllm}, then frozen and hashed,
so every verifier and precision scores exactly the same candidates. Candidate accuracy
is 52.1\%, 88.7\% of candidates have an extractable final answer, and 6.5\% reach the
length limit. A secondary generator, Llama-3.2-3B-Instruct, produces pools for a
300-problem subset (exploratory; Section~\ref{sec:robustness}).

\subsection{Verifiers}\label{sec:verifiers}

The main study uses three outcome reward models (ORMs) from one family,
Skywork-Reward-V2-Qwen3 at 1.7B, 4B and 8B parameters \cite{skywork2025v2}, and one
process reward model (PRM), Qwen2.5-Math-PRM-7B \cite{zhang2025prm}. We chose the Skywork family because it provides several sizes on one backbone, all
trained on the same curated preference data \cite{skywork2025v2}, which lets us compare
size and precision within a family, and because they reported state-of-the-art results on seven reward-model benchmarks at
release. The Skywork models are
general preference reward models, trained on broad human-preference data rather than
on mathematical correctness. We use them as outcome verifiers by scoring the full
(problem, response) pair with their scalar head. The PRM, by contrast, is trained for
mathematical step correctness, so the study covers both kinds of supervision. Its step
rewards are aggregated by their minimum, and the last-step reward is reported as a
sensitivity analysis. To test generality beyond the Skywork-Qwen family, two further
ORMs on a different backbone are evaluated in exploratory experiments:
Skywork-Reward-V2-Llama-3.1-8B (same training-data family, Llama~3.1 backbone) and
Llama-3.1-Tulu-3-8B-RM \cite{lambert2024tulu3} (different laboratory and data). All
models are run with their documented input formats (\ref{app:impl}).

\subsection{Precisions and deployments}\label{sec:precisions}

Quantization is weight-only. The weights of all linear layers inside the transformer
blocks are quantized, while embeddings, norms, the reward head, activations and
computation stay in BF16. ``BF16'' denotes the uncompressed model. The conditions are:
\begin{itemize}
\item \textbf{NF4}: 4-bit NormalFloat \cite{dettmers2023qlora}, both simulated and
deployed with \texttt{bitsandbytes} (bit-identical scores);
\item \textbf{GPTQ-W4} and \textbf{GPTQ-W3}: GPTQ \cite{frantar2023gptq} with group
size 128, calibrated on 128 WikiText-2 sequences of 512 tokens (generic calibration);
\item \textbf{AWQ-W4} \cite{lin2024awq} with the same calibration, and round-to-nearest
\textbf{RTN-W4} and \textbf{RTN-W8} for an exploratory dose-response.
\end{itemize}
Accuracy experiments use simulated (``fake'') quantization, in which weights are
quantized and then dequantized back to BF16. The scores then differ from BF16 only by the
rounding of the weights, so every method and bit width runs on the same kernels and the
comparison is not mixed with kernel-specific numerics. Simulation also covers conditions
that have no packed deployment here (GPTQ-W3, AWQ, RTN and the PRM). We check that it
represents real deployment by comparing it with packed GPTQ-W4 kernels
(Section~\ref{sec:robustness}). For the memory
study, the GPTQ-W4 ORMs are exported and run with packed 4-bit kernels
(\texttt{gptqmodel}, Marlin). The PRM has no packed GPTQ deployment because the export
does not support its remote modelling code, so its deployed 4-bit condition is NF4.
Every quantized model is audited: every target layer must differ from BF16 and every
other parameter must stay bit-identical.

Verifier memory is the peak allocated GPU memory, reported in GiB ($2^{30}$ bytes; 1~GiB
$\approx$ 1.074~GB, where 1~GB = $10^9$ bytes). GB appears only where a source states
it, such as GPU specifications. Memory is measured while
scoring a fixed workload of 64 sequences of 900 tokens on an RTX 4090 with no generator
resident.

\subsection{Selection rules and baselines}\label{sec:selection}

Best-of-$N$ selects the highest-scoring candidate among $N$, and exact score ties are
broken by a fixed random priority (seed 11) shared by all conditions. For $N<32$,
accuracy is averaged over 200 random subsets per problem, drawn identically for every
condition (seed 7); $N=32$ is the full pool. We report two verifier-free baselines.
Majority vote is a plurality vote over canonical answer clusters, formed by numeric
normalisation followed by symbolic equivalence with \texttt{math\_verify}. Candidates
without an extractable answer do not vote, and a pool with no extractable answer counts
as incorrect. The most frequent cluster wins, with ties broken in favour of the cluster
that appears first in candidate order, and the vote is correct if the winning cluster
contains a correct candidate. Majority vote therefore does not depend on the verifier.
The oracle (pass@32) counts a problem as solved if at least one of its 32 candidates is
correct. It is the upper bound for any rule that selects one candidate from the pool.

\subsection{Grading}\label{sec:grading}

The prespecified grader extracts the last \texttt{\textbackslash boxed\{\}} answer, or
failing that a GSM8K-style ``\#\#\#\#'' line or an explicit ``the answer is'' phrase. It
normalises numbers and compares MATH answers symbolically. A strict grader can confound
verifier preferences with answer format, so we add three checks. The first is a lenient
grader, used for sensitivity analyses only, which takes the last number or expression in
the final lines when no answer is extractable. The second is a blinded human audit
(\ref{app:audit}), and the third re-grades GSM8K with the revised GSM8K-Platinum labels
\cite{vendrow2025platinum}. The endpoint throughout is final-answer correctness.

Under the strict grader a winner without an extractable answer always counts as wrong,
so Best-of-$N$ accuracy factorises as $A = p\,a$, where $p$ is the share of winners with
an extractable answer and $a$ the accuracy among them. A change from BF16 (subscript 0)
to a quantized condition (subscript 1) then splits exactly into a format term and a
content term, $\Delta = (p_1 - p_0)\,a_0 + p_1\,(a_1 - a_0)$.

\subsection{Metrics and controls}\label{sec:metrics}

For verifier $v$ and condition $c$, $A_N$ is Best-of-$N$ accuracy and
$\Delta_N = A_N^{c} - A_N^{\mathrm{BF16}}$, computed per problem and averaged. At $N=32$ we also compare each quantized selection with the BF16 selection of the
same problem and report four shares of problems: the selected candidate changes; the
selected canonical answer changes; a correct BF16 selection becomes wrong (harm); and a
wrong BF16 selection becomes correct (benefit). $\Delta_{32}$ equals benefit minus
harm. Two controls measure how much the selection moves without any quantization.
\emph{bf16-null} rescores in BF16 with a different batch layout, which changes padding
and the order of floating-point operations and hence perturbs scores slightly.
\emph{bf16-repeat} rescores with the same layout. The repeat is
bit-identical. The null control changes the selected candidate on 7--11\% of
main-study problems with $|\Delta_{32}| \le 0.3$\pp, and on 10.9\% and 15.0\% (H6) and
5.5\% and 8.7\% (H7) of problems for the ORM and the PRM.

\subsection{Statistics}\label{sec:stats}

All inference uses a paired percentile bootstrap over problems with 10,000 resamples
(seed 7). Problems are resampled with replacement from the full population. Every
quantity computed on a problem (both conditions, all verifiers entering a contrast, the
majority-vote outcome) is resampled jointly, which keeps within-problem correlations.
\begin{itemize}
\item \emph{Differences} are reported with 95\% percentile intervals and two-sided
centred-bootstrap $p$-values.
\item \emph{Equivalence} at the prespecified $\pm3$\pp\ margin is tested by two one-sided
tests (TOST), $p_{\mathrm{TOST}} = \max\{P^*(\bar d \le -0.03),\, P^*(\bar d \ge
0.03)\}$. The test inverts the percentile interval, so $p_{\mathrm{TOST}} \le \alpha$
if and only if the $(1-2\alpha)$ interval lies inside the margin. We checked its size by
simulation with the true mean on either margin, drawing $n=800$ per-problem differences
from $\{-1,-\tfrac12,0,\tfrac12,1\}$ with probabilities $(0.075, 0.03, 0.82, 0.03,
0.045)$ for the lower margin and the mirror image for the upper one (2,000 replicates of
2,000 resamples). The rejection rates were 0.056 and 0.054 at nominal 0.05, 0.024 and
0.025 at 0.025, and 0.006 and 0.004 at 0.0125, with a Monte Carlo standard error of about
0.005. The $\pm3$\pp\ margin was fixed in the plan before any evaluation score existed.
On the c1 population it corresponds to a 13\% relative increase in errors.
\item \emph{Value retention} (H6b, H7b) is tested one-sided, $p = P^*(\bar R \le 0)$
(Section~\ref{sec:h6}).
\item \emph{Resolution.} Bootstrap $p$-values are Monte Carlo estimates with resolution
$10^{-4}$. When no resample crosses the boundary we report the count (0/$10^4$) instead
of $p=0$, and the corresponding 95\% Monte Carlo upper bound is $3\cdot10^{-4}$.
\item \emph{Multiplicity} is controlled by Holm correction within each prespecified
family (c1--c5, H6a/b and H7a/b). The $k$-th smallest of $m$ raw $p$-values is
multiplied by $m-k+1$ and the cumulative maximum is taken. A test of the family that was
not run keeps $p=1$. In the main-study family, c1 enters with the two-sided $p$-value of
its difference test, as implemented when the family was fixed. Its decision rule is
equivalence, judged by whether its 90\% interval lies inside $\pm3$\pp\
(Table~\ref{tab:confirmatory}).
\end{itemize}
\noindent The five main-study contrasts are the following. c1 is the equivalence of
$\Delta_{32}$ for the 8B ORM at GPTQ-W4 on the original 1,200-problem manifest. c2
compares the deployment selected on development data with the best BF16 deployment under
the 6~GiB verifier memory ceiling. c3 compares CH-GPTQ, a GPTQ calibration with a
contrastive Hessian term built from winner--challenger pairs, with an adapted DPQ recipe
\cite{yang2026dpq}; it was not run (\ref{app:deviations}). c4 is $\Delta_{32}$ on the
1,200 MATH problems for the 4B and 8B ORMs at GPTQ-W4, averaged within problem, and c5 is
the change of that loss from $N=4$ to $N=32$.
